\section{AMC 10 10A 2013}
        \begin{enumerate}
        	\item (\textit{AMC 10 10A 2013}) A taxi ride costs $1.50 plus $0.25 per mile traveled.  How much does a 5-mile taxi ride cost?


 $\textbf{(A)}\ 2.25 \qquad\textbf{(B)}\ 2.50  \qquad\textbf{(C)}\ 2.75 \qquad\textbf{(D)}\ 3.00 \qquad\textbf{(E)}\ 3.75$


 
 \textbf{Answer: }C



\textbf{Solution}

There are five miles which need to be traveled. The cost of these five miles is $(0.25\cdot5) = 1.25$.   Adding this to $1.50$, we get $\boxed{\textbf{(C) }2.75}$

 


	\item (\textit{AMC 10 10A 2013}) Alice is making a batch of cookies and needs $2\frac{1}{2}$ cups of sugar.  Unfortunately, her measuring cup holds only $\frac{1}{4}$ cup of sugar.  How many times must she fill that cup to get the correct amount of sugar?


 $\textbf{(A)}\ 8 \qquad\textbf{(B)}\ 10  \qquad\textbf{(C)}\ 12 \qquad\textbf{(D)}\ 16 \qquad\textbf{(E)}\ 20$


 
 \textbf{Answer: }B



\textbf{Solution}

To get how many cups we need, we realize that we simply need to divide the number of cups needed by the number of cups collected in her measuring cup each time.  Thus, we need to evaluate the fraction $\frac{2\frac{1}{2}}{\frac{1}{4}}$.  Simplifying, this is equal to $\frac{5}{2}(4) = \boxed{\textbf{(B) }10}$

 


	\item (\textit{AMC 10 10A 2013}) Square $ABCD$ has side length $10$. Point $E$ is on $\overline{BC}$, and the area of $\triangle ABE$
is $40$. What is $BE$?

                    \begin{figure}[H]
                        \centering
                        \includegraphics[width=0.5\textwidth]{Problem/AMC_10_10A_2013_3_Problem_Figure_1.png}
                    \end{figure}
                    


 $\textbf{(A)}\ 4\qquad\textbf{(B)}\ 5\qquad\textbf{(C)}\ 6\qquad\textbf{(D)}\ 7\qquad\textbf{(E)}\ 8$


 
 \textbf{Answer: }E



\textbf{Solution}

We are given that the area of $\triangle ABE$ is $40$, and that $AB = 10$.

 The area of a triangle is:

 $A = \frac{bh}{2}$

 Using $AB$ as the height of $\triangle ABE$,

 $40 = \frac{10b}{2}$

 and solving for $b$,

 $b = 8$, which is $\boxed{\textbf{(E)}}$

 


	\item (\textit{AMC 10 10A 2013}) A softball team played ten games, scoring 1, 2, 3, 4, 5, 6, 7, 8, 9, and 10 runs.  They lost by one run in exactly five games.  In each of their other games, they scored twice as many runs as their opponent. How many total runs did their opponents score?


 $\textbf{(A)}\ 35 \qquad\textbf{(B)}\ 40 \qquad\textbf{(C)}\ 45 \qquad\textbf{(D)}\ 50 \qquad\textbf{(E)}\ 55$


 
 \textbf{Answer: }C



\textbf{Solution 1}

We know that, for the games where they scored an odd number of runs, they cannot have scored twice as many runs as their opponents, as odd numbers are not divisible by $2$.  Thus, from this, we know that the five games where they lost by one run were when they scored $1$, $3$, $5$, $7$, and $9$ runs, and the others are where they scored twice as many runs.  We can make the following chart:

 $\begin{tabular}{|l|l|} \hline Them & Opponent \\ \hline 1 & 2 \\ 2 & 1 \\ 3 & 4 \\ 4 & 2 \\ 5 & 6 \\ 6 & 3 \\ 7 & 8 \\ 8 & 4 \\ 9 & 10 \\ 10 & 5 \\ \hline \end{tabular}$

 The sum of their opponent's scores is $2 + 1 + 4 + 2 + 6 + 3 + 8 + 4 + 10 + 5 = \boxed{\textbf{(C) }45}$

 



\textbf{Solution 2 (patterns and easier arithmetic)}

The team must've won the games with the even runs and lost the ones with the odd runs. The opponents will have an arithmetic sequence of runs, $1,2,3,4,5$ when the team has even runs. The opponents will have an arithmetic sequence of even runs, $2,4,6,8,10$, when the team has odd runs. The sum of their runs is $3\cdot(1+2+3+4+5)=3\cdot15=45=\fbox{C}$~dragnin

 



\textbf{Solution 3}

They scored $1+2+3+\dots+10=55$. On the five games where they scored even, their opponents must have scored half of their point so the opponent scored $\frac{2+4+6+8+10}{2}=15$ less than them. On the other games they scored $1$ more so the opponent scored 5 more than them in the odd games. Then the opponent's score is $55-15+5=45$.

 


	\item (\textit{AMC 10 10A 2013}) Tom, Dorothy, and Hammy went on a vacation and agreed to split the costs evenly.  During their trip Tom paid $\$105$, Dorothy paid $\$125$, and Hammy paid $\$175$.  In order to share costs equally, Tom gave Hammy $t$ dollars, and Dorothy gave Hammy $d$ dollars.  What is $t-d$?


 $\textbf{(A)}\ 15\qquad\textbf{(B)}\ 20\qquad\textbf{(C)}\ 25\qquad\textbf{(D)}\ 30\qquad\textbf{(E)}\ 35$


 
 \textbf{Answer: }B



\textbf{Solution 1}

The total amount paid is $105 + 125 + 175 = 405$.  To get how much each should have paid, we do $405/3 = 135$.

 Thus, we know that Tom needs to give Sammy 30 dollars, and Dorothy 10 dollars.  This means that $t-d = 30 - 10 = \boxed{\textbf{(B) }20}$.

 



\textbf{Solution 2}

The difference in the money that Tom paid and Dorothy paid is $20$.  In order for them both to have paid the same amount, Tom must pay $20$ more than Dorothy.  The answer is $\boxed{{(B)20}}$.

 



\textbf{Solution 3}

The meaning of sharing costs equally is meaning that, after the vacation, they are equally dividing the money in a way such that, each person would have the same amount left. As each person spends an amount of money, greater than 100, let it be that they all had $200$ dollars to spend. This means that after the vacation we want the amount of money, they currently have. After the trip, Tom would've $95$ dollars, Dorothy would've $75$ dollars, and Hammy had $25$ dollars. This gives us a total of $95+75+25=195$ dollars. 

 We want to equally split this money, as that is what happens after splitting the cost equally. This means that we want Dorothy, Tom, and Hammy to each have $65$ dollars. We know that Tom gave Hammy $t$ dollars meaning that we want to split this money first. As Tom gives money to no one else, we want him to reach $65$ dollars in this trade, meaning that as Tom has $95$ dollars and Hammy has $25$ dollars, we can do a trade of $30$ so $t=30$. After this trade, we get that Tom has $65$ dollars, Hammy has $55$ dollars, and Dorothy has $75$ dollars.

 Next trade is where Dorothy gives $d$ dollars to Hammy. Dorothy has $75$ dollars and Hammy has $55$ dollars. As both of these don't have $65$ dollars and this is the last trade, we need to make sure both have $65$ dollars at the end. This is possible if $d=10$

 We want to find $t-d=30-10=\boxed{B) 20}$

 


	\item (\textit{AMC 10 10A 2013}) Kris and his five brothers are ages 3, 5, 7, 9, 11, and 13.  One afternoon two of his brothers whose ages sum to 16 went to the movies, two brothers younger than 10 went to play baseball, and Kris and the 5-year-old stayed home.  How old is Kris?


 $\textbf{(A)}\ 3 \qquad\textbf{(B)}\ 7  \qquad\textbf{(C)}\ 9 \qquad\textbf{(D)}\ 11 \qquad\textbf{(E)}\ 13$


 
 \textbf{Answer: }D



\textbf{Solution 1}

Because the $5$-year-old stayed home, we know that the $11$-year-old did not go to the movies, as the $5$-year-old did not and $11+5=16$.  Also, the $11$-year-old could not have gone to play baseball, as he is older than $10$.  Thus, the $11$-year-old must have stayed home, so Kris is $\boxed{\textbf{(D) }11}$

 



\textbf{Solution 2}

There are only $4$ kids who are under $10$ but since the $5$-year old stayed home, the only possible ages who went to play baseball are the brothers who are $3,7,9$, either $13+3$ or $7+9$ is $16$ but since we need $2$ kids to go to baseball who are under $10$, $13,3$ must have been the pair to go to the movies and $9,7$ must have went to baseball, so only the $11$-year old is left, which is answer choice $\boxed{\textbf{(D) }11}$

 


	\item (\textit{AMC 10 10A 2013}) A student must choose a program of four courses from a menu of courses consisting of English, Algebra, Geometry, History, Art, and Latin. This program must contain English and at least one mathematics course. In how many ways can this program be chosen?


 $\textbf{(A)}\ 6\qquad\textbf{(B)}\ 8\qquad\textbf{(C)}\ 9\qquad\textbf{(D)}\ 12\qquad\textbf{(E)}\ 16$


 
 \textbf{Answer: }C



\textbf{Solution 1}

Let us split this up into two cases.

 Case $1$: The student chooses both algebra and geometry.

 This means that $3$ courses have already been chosen.  We have $3$ more options for the last course, so there are $3$ possibilities here.  

 Case $2$: The student chooses one or the other.  

 Here, we simply count how many ways we can do one, multiply by $2$, and then add to the previous.  

 Assume the mathematics course is algebra.  This means that we can choose $2$ of History, Art, and Latin, which is simply $\dbinom{3}{2} = 3$.  If it is geometry, we have another $3$ options, so we have a total of $6$ options if only one mathematics course is chosen.

 Thus, overall, we can choose a program in $6 + 3 = \boxed{\textbf{(C) }9}$ ways

 



\textbf{Solution 2}

We can use complementary counting. Since there must be an English class, we will add that to our list of classes leaving $3$ remaining spots for rest of classes. We are also told that there needs to be at least one math class. This calls for complementary counting. The total number of ways of choosing $3$ classes out of the $5$ is $\binom{5}3$. The total number of ways of choosing only non-mathematical classes is $\binom{3}3$. Therefore the amount of ways you can pick classes with at least one math class is $\binom{5}3-\binom{3}3=10-1=\boxed{\textbf{(C) }9}$ ways.

 



\textbf{Solution 3}

Similar to Solution 1, note that for Case 2 of solution the answer is simply $\binom{4}2$ and for the first case it is $\binom{3}2$  hence $\boxed{\textbf{(C) }9}$ ways

 



\textbf{Solution 4}

There are $\binom{4}{2}$ ways to choose two courses given we have chosen English and Algebra; another $\binom{4}{2}$ ways to choose courses given we have English and Geometry. The total, $12$, counts the case with both Algebra and Geometry twice, so we subtract $\binom{3}{1} = 3$ - the number of ways to choose one course from History, Art, and Latin - from $12$ to get $\boxed{ \textbf{ (C) } 9 }$ ways.

 


	\item (\textit{AMC 10 10A 2013}) What is the value of $\frac{2^{2014}+2^{2012}}{2^{2014}-2^{2012}} ?$


 $\textbf{(A)}\ -1 \qquad\textbf{(B)}\ 1  \qquad\textbf{(C)}\ \frac{5}{3} \qquad\textbf{(D)}\ 2013 \qquad\textbf{(E)}\ 2^{4024}$


 
 \textbf{Answer: }C



\textbf{Solution}

Factoring out, we get: $\frac{2^{2012}(2^2 + 1)}{2^{2012}(2^2-1)}$.  

 Cancelling out the $2^{2012}$ from the numerator and denominator, we see that it simplifies to $\boxed{\textbf{(C) }\frac{5}{3}}$.

 



\textbf{Solution 2}

Let $x=2^{2012}$.

 Then the given expression is equal to $\frac{4x+x}{4x-x}=\frac{5x}{3x}=\boxed{\textbf{(C) }\frac{5}{3}}$.

 



\textbf{Solution 3}

We observe that all the answer choices other than $\textbf{(C)}$ are integer.

 Claim: $\frac{2^{2014}+2^{2012}}{2^{2014}-2^{2012}}$ is not an integer.

 Proof:We know that if $b>a$ and $\operatorname{gcd}(a,b) = a$, then $a|b$.

 So all we need to prove is that $\operatorname{gcd}(2^{2014}+2^{2012}, 2^{2014}-2^{2012}) \neq 2^{2014}+2^{2012} \text{ or } 2^{2014}-2^{2012}$. Since the numerator is larger than the denominator, the only possible case is $\gcd(2^{2014}+2^{2012}, 2^{2014}-2^{2012}) = 2^{2014}-2^{2012}$. We have to prove that this is impossible. Using the Euclidean Algorithm we can simplify to obtain\begin{align*}\gcd(2^{2014}+2^{2012}, 2^{2014}-2^{2012}) &= \gcd(2^{2014}+2^{2012}-2^{2014}+2^{2012}, 2^{2014}-2^{2012})\\&= \gcd(2 \cdot 2^{2012}, 2^{2014}-2^{2012})\\&= \gcd(2^{2013}, 2^{2014}-2^{2012})\end{align*}
We notice that $2^{2014}-2^{2012}$ can be factored since $2^{2012}$ is a common factor. Factoring gives us $2^{2014}-2^{2012} = 2^{2012}(2^2-1) = 3 \cdot 2^{2012}$. From this we know that $3$ is a factor of the original expression. Hence $\gcd(2^{2013}, 2^{2014}-2^{2012}) = \gcd(2^{2013}, 3 \cdot 2^{2012}) = 2^{2013} \neq 2^{2012} \cdot 3$. So $2^{2014}-2^{2012} \nmid 2^{2014}+2^{2012}$. So the fraction cannot be simplified into an integer. Hence the answer is the only fraction, $\boxed{\textbf{(C)}}$.

 


	\item (\textit{AMC 10 10A 2013}) In a recent basketball game, Shenille attempted only three-point shots and two-point shots.  She was successful on $20\%$ of her three-point shots and $30\%$ of her two-point shots.  Shenille attempted $30$ shots.  How many points did she score?


 $\textbf{(A)}\ 12 \qquad\textbf{(B)}\ 18  \qquad\textbf{(C)}\ 24 \qquad\textbf{(D)}\ 30 \qquad\textbf{(E)}\ 36$


 
 \textbf{Answer: }B



\textbf{Solution 1}

Let the number of attempted three-point shots be $x$ and the number of attempted two-point shots be $y$.  We know that $x+y=30$, and we need to evaluate $3(0.2x) + 2(0.3y)$, as we know that the three-point shots are worth $3$ points and that she made $20$\% of them and that the two-point shots are worth $2$ and that she made $30$\% of them.

 Simplifying, we see that this is equal to $0.6x + 0.6y = 0.6(x+y)$.  Plugging in $x+y=30$, we get $0.6(30) = \boxed{\textbf{(B) }18}$

 



\textbf{Solution 2}

The problem statement implies that it doesn't matter how many two-point shots or three-point shots are attempted. If we assume that Shenille only attempts three-pointers, then she makes $0.2 \cdot 30 = 6$ shots, which are worth $6 \cdot 3 = \boxed{\textbf{(B) }18}$ points. If we assume Shenille only attempts two-pointers, then she makes $0.3 \cdot 30 = 9$ shots, which are worth $9 \cdot 2 = 18$ points.

 


	\item (\textit{AMC 10 10A 2013}) A flower bouquet contains pink roses, red roses, pink carnations, and red carnations.  One third of the pink flowers are roses, three fourths of the red flowers are carnations, and six tenths of the flowers are pink.  What percent of the flowers are carnations?


 $\textbf{(A)}\ 15 \qquad\textbf{(B)}\ 30  \qquad\textbf{(C)}\ 40 \qquad\textbf{(D)}\ 60 \qquad\textbf{(E)}\ 70$


 
 \textbf{Answer: }E



\textbf{Solution 1}

Let the total amount of flowers be $x$.  Thus, the number of pink flowers is $0.6x$, and the number of red flowers is $0.4x$.  The number of pink carnations is $\frac{2}{3}(0.6x) = 0.4x$ and the number of red carnations is $\frac{3}{4}(0.4x) = 0.3x$.  Summing these, the total number of carnations is $0.4x+0.3x=0.7x$.  Dividing, we see that  $\frac{0.7x}{x} = 0.7 = \boxed{\textbf{(E) }70\%}$

 



\textbf{Solution 2}

We have $\dfrac15$ for pink roses,$1-\dfrac6{10}=\dfrac4{10}=\dfrac25$ red flowers, $\dfrac6{10}-\dfrac15=\dfrac35-\dfrac15=\dfrac25$ pink carnations, $\dfrac25\cdot \dfrac34=\dfrac6{20}=\dfrac3{10}$ red carnations we add them up to get $\dfrac25+\dfrac3{10}=\dfrac7{10}=70\%$ so our final answer is \textbf{70\%} or $\boxed{\textbf{(E) }70\%}$

 


	\item (\textit{AMC 10 10A 2013}) A student council must select a two-person welcoming committee and a three-person planning committee from among its members. There are exactly 10 ways to select a two-person team for the welcoming committee. It is possible for students to serve on both committees. In how many different ways can a three-person planning committee be selected?


 $\textbf{(A)}\ 10\qquad\textbf{(B)}\ 12\qquad\textbf{(C)}\ 15\qquad\textbf{(D)}\ 18\qquad\textbf{(E)}\ 25$


 
 \textbf{Answer: }A



\textbf{Solution 1}

Let the number of students on the council be $x$.  To select a two-person committee, we can select a "first person" and a "second person." There are $x$ choices to select a first person; subsequently, there are $x-1$ choices for the second person. This gives a preliminary count of $x(x-1)$ ways to choose a two-person committee. However, this accounts for the order of committees. To understand this, suppose that Alice and Bob are two students in the council. If we choose Alice and then Bob, that is the same as choosing Bob and then Alice and so latter and former arrangements would be considered the same. Therefore, we have to divide by $2$ to account for overcounting. Thus, there are $\dfrac{x(x-1)} 2=10$ ways to choose the two-person committee. Solving this equation, we find that $5$ and $-4$ are integer solutions. $-4$ is a ridiculous situation, so there are $5$ people on the student council. The solution is $\dbinom 5 3=10\implies \boxed{\textbf{A}}$.

 



\textbf{Solution 2 (faster)}

To choose $2$ people from $n$ total people and get $10$ as a result, we can establish the equation $\binom{n}{2}=10$ which we can easily see $n=5$, so there are $5$ people. The question asks how many ways to choose $3$ people from the $5$, so there are $\binom{5}{3}=\boxed{\textbf{(A)}10}$ ways. (Note that Solution 1 is more rigorous, making sure that $n = 5$ is the only possible $n$.)

 


	\item (\textit{AMC 10 10A 2013}) In $\triangle ABC$, $AB=AC=28$ and $BC=20$.  Points $D,E,$ and $F$ are on sides $\overline{AB}$, $\overline{BC}$, and $\overline{AC}$, respectively, such that $\overline{DE}$ and $\overline{EF}$ are parallel to $\overline{AC}$ and $\overline{AB}$, respectively.  What is the perimeter of parallelogram $ADEF$?


 
                    \begin{figure}[H]
                        \centering
                        \includegraphics[width=0.5\textwidth]{Problem/AMC_10_10A_2013_12_Problem_Figure_1.png}
                    \end{figure}
                    


 $\textbf{(A) }48\qquad \textbf{(B) }52\qquad \textbf{(C) }56\qquad \textbf{(D) }60\qquad \textbf{(E) }72\qquad$


 
 \textbf{Answer: }C



\textbf{Solution 1}

Note that because $\overline{DE}$ and $\overline{EF}$ are parallel to the sides of $\triangle ABC$, the internal triangles $\triangle BDE$ and $\triangle EFC$ are similar to $\triangle ABC$, and are therefore also isosceles triangles.

 It follows that $BD = DE$. Thus, $AD + DE = AD + DB = AB = 28$.

 The opposite sides of parallelograms are equal (you can prove this fact simply by drawing the diagonal of the parallelogram and proving that the two resulting triangles are congruent by SSS), so the perimeter is $2 \times (AD + DE) = 56\implies \boxed{\textbf{(C)}}$.

 



\textbf{Solution 2}

We will be using logic to solve this problem :-)! 

 The ratio \( AB/BC = DB/BE \) is just \( 28/20 = DB/BE \) which implies, through cross multiplication, that \( 28BE = 20DB \). Similarly, the ratio \( AB/BC = FE/EC \) is just \( 28/20 = FE/EC \), and thus, also implies that \( 28EC = 20FE \). 

 For Reference, we shall also put \( DB/BE = FE/EC \), or \( DB \cdot EC = FE \cdot BE \). 

 Analyzing \( 28BE=20DB \), we find out that the number of ordered pairs \( (BE,DB) \) that work are \( (5, 7) \), \( (10, 14) \) and \( (15, 21) \). 

 The most logical choice for BE is \( BE = 15 \), \( DB = 21 \). 

 Because \( BE + EC = 20 \), We know that \( EC = 5 \), and thus, by using the same logic, \( FE = 7 \). 

 The perimeter of parallelogram \( ADEF \) is none other than \( 2(21) + 2(7) = 42 + 14 = \)$\boxed{56 \Rightarrow \textbf{(C)}}$ 

 



\textbf{Solution 3}

Similar to Solution 1, we find that the internal triangles $\triangle BDE$ and $\triangle EFC$ are similar to $\triangle ABC$, and are therefore also isosceles triangles and similar to each other. We know that the perimeter of $ADEF$ can be found by $2(l + w)$, and $DE$ will be the length and the width will be $EF$. These can be found by looking at the long sides of $\triangle DBE$ and $\triangle FEC$ respectively. We then find the ratio of long sides to the short side of $\triangle ABC$ to be $7/5$ by $28/20$, which applies to the other triangles since they are similar. We set up an expression, calling $BE$ as $x$ and $EC$ as $20 - x$, and substitute it into the perimeter equation knowing our long sides from the ratio:

 $2(7x/5 + 7/5(20-x))$

 $2(7x/5 + 28 - 7x/5)$

 $2(28)$

 $56\implies \boxed{\textbf{(C)}}$ 

 



\textbf{Solution 4}

The problem implies that where you choose $D$, $E$, and $F$ are irrelevant. This means that to make it easier, choose it such that $D$, $E$, and $F$ are the midpoints of $\overline{AB}$, $\overline{BC}$, and $\overline{AC}$ respectively. $28 \div 2 = 14$, $14 \times 4 = 56$, so $\boxed{\textbf{(C)}}$. 

 



\textbf{Solution 5}

Drawing the diagram with a ruler and compass (and scaling back by x4), we can draw approximate parallel lines. This yields about 14, but we need to multiply by 4 to get 56, or $\boxed{(C)}$.

 


	\item (\textit{AMC 10 10A 2013}) How many three-digit numbers are not divisible by $5$, have digits that sum to less than $20$, and have the first digit equal to the third digit?


 $\textbf{(A)}\ 52 \qquad\textbf{(B)}\ 60  \qquad\textbf{(C)}\ 66 \qquad\textbf{(D)}\ 68 \qquad\textbf{(E)}\ 70$


 
 \textbf{Answer: }B



\textbf{Solution}

We use a casework approach to solve the problem. These three digit numbers are of the form $\overline{xyx}$.($\overline{abc}$ denotes the number $100a+10b+c$).  We see that $x\neq 0$ and $x\neq 5$, as $x=0$ does not yield a three-digit integer and $x=5$ yields a number divisible by 5.

 The second condition is that the sum $2x+y<20$.  When $x$ is $1$, $2$, $3$, or $4$, $y$ can be any digit from $0$ to $9$, as $2x<10$.  This yields $10(4) = 40$ numbers.  

 When $x=6$, we see that $12+y<20$ so $y<8$.  This yields $8$ more numbers.  

 When $x=7$,  $14+y<20$ so $y<6$.  This yields $6$ more numbers.  

 When $x=8$, $16+y<20$ so $y<4$.  This yields $4$ more numbers.  

 When $x=9$, $18+y<20$ so $y<2$.  This yields $2$ more numbers.  

 Summing, we get $40 + 8 + 6 + 4 + 2 = \boxed{\textbf{(B) }60}$

 


	\item (\textit{AMC 10 10A 2013}) A solid cube of side length $1$ is removed from each corner of a solid cube of side length $3$.  How many edges does the remaining solid have?


 $\textbf{(A) }36\qquad\textbf{(B) }60\qquad\textbf{(C) }72\qquad\textbf{(D) }84\qquad\textbf{(E) }108\qquad$


 
 \textbf{Answer: }D



\textbf{Solution 1}

We can use Euler's polyhedron formula that says that $F+V=E+2$.  We know that there are originally $6$ faces on the cube, and each corner cube creates $3$ more.  $6+8(3) = 30$.  In addition, each cube creates $7$ new vertices while taking away the original $8$, yielding $8(7) = 56$ vertices.  Thus $E+2=56+30$, so $E=\boxed{\textbf{(D) }84}$

 



\textbf{Solution 2}

The removal of each cube adds nine additional edges to the solid. Since a cube initially has $12$ edges and there are eight vertices, the number of edges will be $12 + 9 \times 8 = \boxed{\textbf{(D) } 84}$.

 


	\item (\textit{AMC 10 10A 2013}) Two sides of a triangle have lengths $10$ and $15$.  The length of the altitude to the third side is the average of the lengths of the altitudes to the two given sides.  How long is the third side?


 $\textbf{(A)}\ 6 \qquad\textbf{(B)}\ 8 \qquad\textbf{(C)}\ 9 \qquad\textbf{(D)}\ 12 \qquad\textbf{(E)}\ 18$


 
 \textbf{Answer: }D



\textbf{Solution 1 (Process of Elimination)}

The shortest side length has the longest altitude perpendicular to it. The average of the two altitudes given will be between the lengths of the two altitudes,therefore the length of the side perpendicular to that altitude will be between $10$ and $15$. The only answer choice that meets this requirement is $\boxed{\textbf{(D) }12}$.\textbf{}

 



\textbf{Solution 2}

\textbf{}Let the height to the side of length $15$ be $h_{1}$, the height to the side of length 10 be $h_{2}$, the area be $A$, and the height to the unknown side be $h_{3}$.

 Because the area of a triangle is $\frac{bh}{2}$, we get that $15(h_{1}) = 2A$ and $10(h_{2}) = 2A$, so, setting them equal, $h_{2} = \frac{3h_{1}}{2}$.  From the problem, we know that $2h_{3} = h_{1} + h_{2}$.  Substituting, we get that
 \[h_{3} = 1.25h_{1}.\]  Thus, the side length is going to be $\frac{2A}{1.25h_{1}} = \frac{15}{\frac{5}{4}} = \boxed{\textbf{(D) }12}$.

 



\textbf{Solution 3 (Answer Choices)}

Let $x$ be the height of triangle when the base is $10$ and $y$ is the height of the triangles when the base is $15$. This means the height for when the triangles has the $3$rd side length, the height would be $\dfrac{(x+y)}{2}$ giving us the following equation:

 $\dfrac{10x}{2}=\dfrac{15y}{2}=\dfrac{n(x+y)}{2}$For $n$, we can plug in the answer choices and check when the ratio of $\dfrac{x}{y}$ for $\dfrac{10x}{2}=\dfrac{n(x+y)}{2}$ and $\dfrac{15y}{2}=\dfrac{n(x+y)}{2}$ is the same because we can observe that only for one value, the ratio remains constant. From brute force we get that $\boxed{n=12}$ meaning that $12$ is the $3$rd base.~ Batmanstark

 


	\item (\textit{AMC 10 10A 2013}) A triangle with vertices $(6, 5)$, $(8, -3)$, and $(9, 1)$ is reflected about the line $x=8$ to create a second triangle.  What is the area of the union of the two triangles?


 $\textbf{(A)}\ 9 \qquad\textbf{(B)}\ \frac{28}{3} \qquad\textbf{(C)}\ 10 \qquad\textbf{(D)}\ \frac{31}{3} \qquad\textbf{(E)}\ \frac{32}{3}$


 
 \textbf{Answer: }E



\textbf{Solution 1}

Let $A$ be at $(6, 5)$, B be at $(8, -3)$, and $C$ be at $(9, 1)$.  Reflecting over the line $x=8$, we see that $A' = D =  (10,5)$, $B' = B$ (as the x-coordinate of B is 8), and $C' = E =  (7, 1)$.  Line $AB$ can be represented as $y=-4x+29$, so we see that $E$ is on line $AB$.  

 
                    \begin{figure}[H]
                        \centering
                        \includegraphics[width=0.5\textwidth]{Solution/AMC_10_10A_2013_16_Solution_Figure_1.png}
                    \end{figure}
                    

 
We see that if we connect $A$ to $D$, we get a line of length $4$ (between $(6, 5)$ and $(10,5)$).  The area of $\triangle ABD$ is equal to $\frac{bh}{2} = \frac{4(8)}{2} = 16$.

 Now, let the point of intersection between $AC$ and $DE$ be $F$.  If we can just find the area of $\triangle ADF$ and subtract it from $16$, we are done.

 We realize that because the diagram is symmetric over $x = 8$, the intersection of lines $AC$ and $DE$ should intersect at an x-coordinate of $8$.  We know that the slope of $DE$ is $\frac{5-1}{10-7} = \frac{4}{3}$.  Thus, we can represent the line going through $E$ and $D$ as $y - 1=\frac{4}{3}(x - 7)$.  Plugging in $x = 8$, we find that the y-coordinate of F is $\frac{7}{3}$.  Thus, the height of $\triangle ADF$ is $5 - \frac{7}{3} = \frac{8}{3}$.  Using the formula for the area of a triangle, the area of $\triangle ADF$ is $\frac{16}{3}$.  

 To get our final answer, we must subtract this from $16$.  $[ABD] - [ADF] = 16 - \frac{16}{3} = \boxed{\textbf{(E) }\frac{32}{3}}$

 



\textbf{Solution 2}

First, realize that $E$ is the midpoint of $AB$ and $C$ is the midpoint of $BD$. Connect $A$ to $D$ to form $\triangle ABD$. Let the midpoint of $AD$ be $G$. Connect $B$ to $G$. $BG$ is a median of $\triangle ABD$. 

 
Because $\triangle ABD$ is isosceles, $BG$ is also an altitude of $\triangle ABD$. We know the length of $AD$ and $BG$ from the given coordinates. The area of $\triangle ABD$ is $\frac{bh}{2} = \frac{4(8)}{2} = 16$. 

 
Let the intesection of $AC$, $DE$ and $BG$ be $F$. $F$ is the centroid of $\triangle ABD$. Therefore, it splits $BG$ into $BF={2 \over 3}(BG)$ and $FG={1\over 3}(BG)$. The area of quadrilateral $ABDF = 16\cdot {2 \over 3} = \boxed{\textbf{(E) }\frac{32}{3}}$

 



\textbf{Solution 3}

Since we have to find the area on a coordinate plane, we can use the \href{https://artofproblemsolving.com/wiki/index.php?title=Shoelace\_Theorem}{Shoelace Theorem} to find the area of the intersection. When you reflect it, it makes a quadrilateral. 

 
                    \begin{figure}[H]
                        \centering
                        \includegraphics[width=0.5\textwidth]{Solution/AMC_10_10A_2013_16_Solution_Figure_1.png}
                    \end{figure}
                    

 Since it is reflected around $x=8$, the point $(8,-3)$ remains the same on both. The top right corners are $(6,5)$, and its reflection $(10,5)$. Now to find the 4th point, point F, we can use the equation of the line DE($\frac{4}{3}x - \frac{25}{3}$, and substitute $x=8$, to get $\frac{7}{3}$. Now we can use the theorem: $\frac{1}{2}(((6*-3)+(10*5)+(8*5)+(\frac{7}{3}*8))-((8*5)+(6*5)+(10*\frac{7}{3}) + (8* -3))) = \boxed{\textbf{(E) }\frac{32}{3}}$

 


	\item (\textit{AMC 10 10A 2013}) Daphnia is visited periodically by her three best friends: Alice, Beatrix, and Claire.  Alice visits every third day, Beatrix visits every fourth day, and Claire visits every fifth day.  All three friends visited Daphnia yesterday.  How many days of the next 365-day period will exactly two friends visit her?


 $\textbf{(A)}\ 48\qquad\textbf{(B)}\ 54\qquad\textbf{(C)}\ 60\qquad\textbf{(D)}\ 66\qquad\textbf{(E)}\ 72$


 
 \textbf{Answer: }B



\textbf{Solution 1}

The $365$-day time period can be split up into $6$, $60$-day time periods, because after $60$ days, all three of them visit again (Least common multiple of $3$, $4$, and $5$).You can find how many times each pair of visitors can meet by finding the LCM of their visiting days and dividing that number by 60.Remember to subtract $1$, because you do not wish to count the $60$th day, when all three visit.

 A and B visit $\frac{60}{3 \cdot 4}-1 = 4$ times.

 A and C visit $\frac{60}{3 \cdot 5}-1 = 3$ times.

 B and C visit $\frac{60}{4 \cdot 5}-1 = 2$ times.

 
This is a total of $9$ visits per $60$ day period.Therefore, the total number of $2$-person visits is $9 \cdot 6 = \boxed{\textbf{(B) }54}$.

 
                \begin{itemize}
                    \item Note: We do not have to worry about the numbers over 360: ($361,362,363,364,365$) having 2 factors. This is because we can rewrite
                \end{itemize}


                

 


	\item (\textit{AMC 10 10A 2013}) Let points $A = (0, 0)$, $B = (1, 2)$, $C=(3, 3)$, and $D = (4, 0)$.  Quadrilateral $ABCD$ is cut into equal area pieces by a line passing through $A$.  This line intersects $\overline{CD}$ at point $\bigg(\frac{p}{q}, \frac{r}{s}\bigg)$, where these fractions are in lowest terms.  What is $p+q+r+s$?


 $\textbf{(A)}\ 54\qquad\textbf{(B)}\ 58\qquad\textbf{(C)}\ 62\qquad\textbf{(D)}\ 70\qquad\textbf{(E)}\ 75$


 
 \textbf{Answer: }B



\textbf{Solution 1}

First, we shall find the area of quadrilateral $ABCD$. This can be done in any of three ways:

 \href{https://artofproblemsolving.com/wiki/index.php?title=Pick\%27s\_Theorem}{Pick's Theorem}: $[ABCD] = I + \dfrac{B}{2} - 1 = 5 + \dfrac{7}{2} - 1 = \dfrac{15}{2}.$

 Splitting: Drop perpendiculars from $B$ and $C$ to the x-axis to divide the quadrilateral into triangles and trapezoids, and so the area is $1 + 5 + \dfrac{3}{2} = \dfrac{15}{2}.$

 \href{https://artofproblemsolving.com/wiki/index.php?title=Shoelace\_Theorem}{Shoelace Theorem}: The area is half of $|1 \cdot 3 - 2 \cdot 3 - 3 \cdot 4| = 15$, or $\dfrac{15}{2}$.

 $[ABCD] = \frac{15}{2}$.  Therefore, each equal piece that the line separates $ABCD$ into must have an area of $\frac{15}{4}$.

 Call the point where the line through $A$ intersects $\overline{CD}$ $E$.  We know that $[ADE] = \frac{15}{4} = \frac{bh}{2}$.  Furthermore, we know that $b = 4$, as $AD = 4$.  Thus, solving for $h$, we find that $2h = \frac{15}{4}$, so $h = \frac{15}{8}$.  This gives that the y coordinate of E is $\frac{15}{8}$.  

 Line CD can be expressed as $y = -3x+12$, so the $x$ coordinate of E satisfies $\frac{15}{8} = -3x + 12$.  Solving for $x$, we find that $x = \frac{27}{8}$.

 From this, we know that $E = \left(\frac{27}{8}, \frac{15}{8}\right)$.  $27 + 15 + 8 + 8 = \boxed{\textbf{(B) }58}$

 



\textbf{Solution 2}

Let the point where the altitude from $E$ to $\overline{AD}$ be labeled $F$.Following the steps above, you can find that the height of $\triangle ADE$ is $\frac{15}{8}$, and from there split the base into two parts, $x$, and $4-x$, such that $x$ is the segment from the origin to the point $F$, and $4-x$ is the segment from point $F$ to point $D$. Then, by the \href{https://artofproblemsolving.com/wiki/index.php?title=Pythagorean\_Theorem}{Pythagorean Theorem}, $x=\frac{27}{8}$, and the answer is $\boxed{\textbf{(B) }58}$

 


	\item (\textit{AMC 10 10A 2013}) In base $10$, the number $2013$ ends in the digit $3$.  In base $9$, on the other hand, the same number is written as $(2676)_{9}$ and ends in the digit $6$.  For how many positive integers $b$ does the base-$b$-representation of $2013$ end in the digit $3$?


 $\textbf{(A)}\ 6\qquad\textbf{(B)}\ 9\qquad\textbf{(C)}\ 13\qquad\textbf{(D)}\ 16\qquad\textbf{(E)}\ 18$


 
 \textbf{Answer: }C



\textbf{Solution 1}

We want the integers $b$ such that $2013\equiv 3\pmod{b} \Rightarrow b$ is a factor of $2010$. Since $2010=2 \cdot 3 \cdot 5 \cdot 67$, it has $(1+1)(1+1)(1+1)(1+1)=2^4=16$ factors. Since $b$ cannot equal $1, 2,$ or $3$, as these can not have the digit $3$ in their base representations, our answer is $16-3=\boxed{\textbf{(C) }13}$

 


	\item (\textit{AMC 10 10A 2013}) A unit square is rotated $45^\circ$ about its center.  What is the area of the region swept out by the interior of the square?


 $\textbf{(A)}\ 1 - \frac{\sqrt2}{2} + \frac{\pi}{4}\qquad\textbf{(B)}\ \frac{1}{2} + \frac{\pi}{4} \qquad\textbf{(C)}\ 2 - \sqrt2 + \frac{\pi}{4}$


 $\textbf{(D)}\ \frac{\sqrt2}{2} + \frac{\pi}{4} \qquad\textbf{(E)}\ 1 + \frac{\sqrt2}{4} + \frac{\pi}{8}$


 
 \textbf{Answer: }C



\textbf{Solution 1}

First, we need to see what this looks like.  Below is a diagram.

 
                    \begin{figure}[H]
                        \centering
                        \includegraphics[width=0.5\textwidth]{Solution/AMC_10_10A_2013_20_Solution_Figure_1.png}
                    \end{figure}
                    

 For this square with side length 1, the distance from center to vertex is $r = \frac{\sqrt{2}}{2}$, hence the area is composed of a semicircle of radius $r$, plus $4$ times a parallelogram (or a kite with diagonals of $(\sqrt{2}-1)$ and $r \text{ or} \frac{\sqrt{2}}{2}$) with height $\frac{1}{2}$ and base $\frac{\sqrt{2}}{2(1+\sqrt{2})}$. That is to say, the total area is $\frac{1}{2} \pi \left(\frac{\sqrt{2}}{2}\right)^2 + 4 \frac{\sqrt{2}}{4(1+\sqrt{2})} = \boxed{\textbf{(C) } 2 - \sqrt{2} + \frac{\pi}{4}}$.

 
                    \begin{figure}[H]
                        \centering
                        \includegraphics[width=0.5\textwidth]{Solution/AMC_10_10A_2013_20_Solution_Figure_2.png}
                    \end{figure}
                    (To turn each dart-shaped piece into a parallelogram, cut along the dashed line and flip over one half.)
                    \begin{figure}[H]
                        \centering
                        \includegraphics[width=0.5\textwidth]{Solution/AMC_10_10A_2013_20_Solution_Figure_3.png}
                    \end{figure}
                    Alternatively, you can move the dart-shaped piece to the other side and make a kite.
                    \begin{figure}[H]
                        \centering
                        \includegraphics[width=0.5\textwidth]{Solution/AMC_10_10A_2013_20_Solution_Figure_4.png}
                    \end{figure}
                    

 



\textbf{Solution 2}


                    \begin{figure}[H]
                        \centering
                        \includegraphics[width=0.5\textwidth]{Solution/AMC_10_10A_2013_20_Solution_Figure_5.png}
                    \end{figure}
                    

 $\textbf{(high res image; no labels)}$

 
                    \begin{figure}[H]
                        \centering
                        \includegraphics[width=0.5\textwidth]{Solution/AMC_10_10A_2013_20_Solution_Figure_6.jpg}
                    \end{figure}
                    

 Let $O$ be the center of the square and $C$ be the intersection of $OB$ and $AD$. The desired area consists of the unit square, plus $4$ regions congruent to the region bounded by arc $AB$, $\overline{AC}$, and $\overline{BC}$, plus $4$ triangular regions congruent to right triangle $BCD$. The area of the region bounded by arc $AB$, $\overline{AC}$, and $\overline{BC}$ is $\frac{\text{Area of Circle}-\text{Area of Square}}{8}$. Since the circle has radius $\dfrac{1}{\sqrt {2}}$, the area of the region is $\dfrac{\dfrac{\pi}{2}-1}{8}$, so 4 times the area of that region is $\dfrac{\pi}{4}-\dfrac{1}{2}$. Now we find the area of $\triangle BCD$. $BC=BO-OC=\dfrac{\sqrt {2}}{2}-\dfrac{1}{2}$. Since $\triangle BCD$ is a $45-45-90$ right triangle, the area of $\triangle BCD$ is $\dfrac{BC^2}{2}=\dfrac {\left (\dfrac {\sqrt {2}}{2}-\dfrac{1}{2} \right)^2}{2}$, so $4$ times the area of $\triangle BCD$ is $\dfrac{3}{2}-\sqrt {2}$. Finally, the area of the whole region is $1+ \left(\dfrac {3}{2}-\sqrt {2} \right) + \left(\dfrac{\pi}{4}-\dfrac{1}{2} \right)=\dfrac{\pi}{4}+2-\sqrt {2}$, which we can rewrite as $\boxed{\textbf{(C) }    2 - \sqrt{2} + \frac{\pi}{4}}$.

 


	\item (\textit{AMC 10 10A 2013}) A group of $12$ pirates agree to divide a treasure chest of gold coins among themselves as follows.  The $k^{\text{th}}$ pirate to take a share takes $\frac{k}{12}$ of the coins that remain in the chest.  The number of coins initially in the chest is the smallest number for which this arrangement will allow each pirate to receive a positive whole number of coins.  How many coins does the $12^{\text{th}}$ pirate receive?


 $\textbf{(A)}\ 720\qquad\textbf{(B)}\ 1296\qquad\textbf{(C)}\ 1728\qquad\textbf{(D)}\ 1925\qquad\textbf{(E)}\ 3850$


 
 \textbf{Answer: }D



\textbf{Solution 1}

Let $x$ be the number of coins.  After the $k^{\text{th}}$ pirate takes his share, $\frac{12-k}{12}$ of the original amount is left.  Thus, we know that

 $x \cdot \frac{11}{12} \cdot \frac{10}{12} \cdot \frac{9}{12} \cdot \frac{8}{12} \cdot \frac{7}{12} \cdot \frac{6}{12} \cdot \frac{5}{12} \cdot \frac{4}{12} \cdot \frac{3}{12} \cdot \frac{2}{12} \cdot \frac{1}{12}$ must be an integer.  Simplifying, we get

 
$x \cdot \frac{11}{12} \cdot \frac{5}{6} \cdot \frac{1}{2}  \cdot \frac{7}{12} \cdot \frac{1}{2} \cdot \frac{5}{12} \cdot \frac{1}{3} \cdot \frac{1}{4} \cdot \frac{1}{6} \cdot \frac{1}{12}$.  Now, the minimal $x$ is the denominator of this fraction multiplied out, obviously.  We mentioned before that this product must be an integer.  Specifically, it is an integer and it is the amount that the $12^{\text{th}}$ pirate receives, as he receives $\frac{12}{12} = 1 =$ all of what is remaining.  

 Thus, we know the denominator is canceled out, so the number of gold coins received is going to be the product of the numerators, $11 \cdot 5 \cdot 7 \cdot 5 = \boxed{\textbf{(D) }1925}$.

 



\textbf{Solution 2 (Using the answer choices)}

Solution $1$ mentioned the expression $x \cdot \frac{11}{12} \cdot \frac{10}{12} \cdot ... \cdot \frac{1}{12}$. Note that this is equivalent to $\frac{x \cdot 11!}{12^{11}}$. 

 
We can compute the amount of factors of $2$, $3$, $5$, etc. but this is not necessary. To minimize this expression, we must take out factors of $2$ and $3$, since $12^{11}=2^{22} \cdot 3^{11}$. $11!$ has neither $22$ factors of $2$, nor $11$ factors of $3$. This means that if $11!$ contains $a$ factors of $2$, then $x$ will contain $22-a$ factors of $2$. This also holds for factors of $3$.

 
Thus, once simplified, the expression will have no factors of $2$. It will also have no factors of $3$.

 
Looking at the answer choices, there is only one answer which is not even, which is $\boxed{\textbf{(D) }1925}$.

 



\textbf{Solution 3}

We know that the 11th pirate takes $\frac{11}{12}$ of what is left from the 10th pirate, so we have the 12th pirate taking 
 \[\frac{12}{12}\cdot(1-\frac{11}{12})=1\cdot\frac{1}{12}\] of what is left from the 10th pirate. Similarly, since the 10th pirate takes $\frac{10}{12}$ of what is left from the 9th pirate, we have the 11th pirate taking 
 \[\frac{11}{12}\cdot(1-\frac{10}{12})=\frac{11}{12}\cdot\frac{2}{12}\] of what is left from the 9th pirate. Thus, the 12th pirate takes 
 \[1\cdot\frac{1}{12}\cdot\frac{2}{12}\] of what is left from the 9th pirate. Repeating the method, we can find that the 12th pirate takes 
 \[1\cdot\frac{1}{12}\cdot\frac{2}{12}\cdot\frac{3}{12}\cdot...\cdot\frac{10}{12}\cdot\frac{11}{12}\] of what is left from the 1st pirate, or 
 \[1\cdot\frac{1}{12}\cdot\frac{2}{12}\cdot\frac{3}{12}\cdot...\cdot\frac{10}{12}\cdot\frac{11}{12}\cdot\frac{12}{12}\] of the total amount of coins. 

 Now canceling out the denominator with the numerator as possible, we are left with $\frac{1\cdot5\cdot7\cdot5\cdot11}{...}=\frac{1925}{...}$ with some factors of 12 in the denominator. For this fraction to be an integer, the smallest possible number of coins is the same as the denominator, so the numerator is the number of coins taken by the 12th pirate, or  $1925 \,\boxed{\mathrm{\textbf{(D)}}}$

 


	\item (\textit{AMC 10 10A 2013}) Six spheres of radius $1$ are positioned so that their centers are at the vertices of a regular hexagon of side length $2$. The six spheres are internally tangent to a larger sphere whose center is the center of the hexagon.  An eighth sphere is externally tangent to the six smaller spheres and internally tangent to the larger sphere.  What is the radius of this eighth sphere?


 $\textbf{(A)}\ \sqrt2\qquad\textbf{(B)}\ \frac{3}{2}\qquad\textbf{(C)}\ \frac{5}{3}\qquad\textbf{(D)}\ \sqrt3\qquad\textbf{(E)}\ 2$


 
 \textbf{Answer: }B



\textbf{Solution 1}

Set up an isosceles triangle between the center of the $8$th sphere and two opposite ends of the hexagon. Then set up another triangle between the point of tangency of the $7$th and $8$th spheres, and the points of tangency between the $7$th sphere and $2$ of the original spheres on opposite sides of the hexagon. Express each side length of the triangles in terms of $r$ (the radius of sphere $8$) and $h$ (the height of the first triangle). You can then use Pythagorean Theorem to set up two equations for the two triangles, and find the values of $h$ and $r$.

 $(1+r)^2=2^2+h^2$

 $(3\sqrt{2})^2=3^2+(h+r)^2$

 $r = \boxed{\textbf{(B) }\frac{3}{2}}$

 Note: We can easily tell that h = 3-r because the line connecting the centers of the two opposite spheres is a diameter of the largest sphere and so intersects the center of the largest sphere, so the middle of this line is the center of the largest sphere. From this we know that the center of the line connected to the top of the 8th sphere is 3 and the height therefore is 3-r~Srag

 



\textbf{Solution 2}

We have a regular hexagon with side length $2$ and six spheres on each vertex with radius $1$ that are internally tangent, therefore, drawing radii to the tangent points would create this regular hexagon.

 Imagine a 2D overhead view.  There is a larger sphere which the $6$ spheres are internally tangent to, with the center in the center of the hexagon. To find the radius of the larger sphere we must first, either by prior knowledge or by deducing from the angle sum that the hexagon can be split into $6$ equilateral triangles from its vertices, that the radius is two more than the small radius, or $3$.  

 

                    \begin{figure}[H]
                        \centering
                        \includegraphics[width=0.5\textwidth]{Solution/AMC_10_10A_2013_22_Solution_Figure_1.png}
                    \end{figure}
                    

 diagram made by erics118

 
Now imagine the figure in $3$ dimensions.  The 8th sphere must be sitting atop of the $6$ spheres, which is the only possibility for it to be tangent to all the $6$ small spheres externally and the larger sphere internally. The ring of $6$ small spheres is symmetrical and the 8th sphere will be resting atop it with its center aligned with the diameter of the large sphere.

 We can now create a right triangle with the legs being the line from a vertex of the hexagon to the center of the hexagon and the line from the center of the 7th sphere to the center of the 8th sphere.  Let the radius of our 8th sphere be $r$.  As previously mentioned, the distance from the center of the hexagon to one of its vertices is $2$.  The distance between the centers will be $3-r$.  The hypotenuse will be $r+1$.  

 We now have a right triangle.  Applying the Pythagorean Theorem, $2^2+(3-r)^2=(1+r)^2$.  Expanding and solving for $r$, we find $r=\frac{12}{8}=\boxed{\textbf{(B)}\frac{3}{2}}$.

 
3D Diagram:\href{https://www.desmos.com/3d/nhwgeghuby}{https://www.desmos.com/3d/nhwgeghuby} (turn on translucent surfaces in more options)

 


	\item (\textit{AMC 10 10A 2013}) In $\triangle ABC$, $AB = 86$, and $AC=97$.  A circle with center $A$ and radius $AB$ intersects $\overline{BC}$ at points $B$ and $X$.  Moreover $\overline{BX}$ and $\overline{CX}$ have integer lengths.  What is $BC$?


 $\textbf{(A)}\ 11\qquad\textbf{(B)}\ 28\qquad\textbf{(C)}\ 33\qquad\textbf{(D)}\ 61\qquad\textbf{(E)}\ 72$


 
 \textbf{Answer: }D



\textbf{Solution 1 (Number Theoretic Power of a Point)}

Let $BX = q$, $CX = p$, and $AC$ meets the circle at $Y$ and $Z$, with $Y$ on $AC$.  Then $AZ = AY = 86$.  Using the Power of a Point (Secant-Secant Power Theorem), we get that $p(p+q) = 11(183) = 11 * 3 * 61$.  We know that $p+q>p$, so $p$ is either $3$, $11$, or $33$.  We also know that $p>11$ by the Triangle Inequality on $\triangle ACX$.  Thus, $p$ is $33$ so we get that $BC = p+q = \boxed{\textbf{(D) }61}$.

 



\textbf{Solution 2}

Let $CX=x,$ and $BX=y.$ Let the circle intersect $AC$ at $D$ and the diameter including $AD$ intersect the circle again at $E.$Use power of a point on point C to the circle centered at A.

 So $CX \cdot CB=CD \cdot CE \Rightarrow x(x+y)=(97-86)(97+86) \Rightarrow x(x+y)=3*11*61.$

 Obviously $x+y>x,$ so we have three solution pairs for $(x,x+y)=(1,2013),(3,671),(11,183),(33,61).$By the Triangle Inequality, only $x+y=61$ yields a possible length of $BX+CX=BC$.

 Therefore, the answer is  $\boxed{\textbf{(D) }61}.$

 



\textbf{Solution 3}


                    \begin{figure}[H]
                        \centering
                        \includegraphics[width=0.5\textwidth]{Solution/AMC_10_10A_2013_23_Solution_Figure_1.png}
                    \end{figure}
                    

 We first draw the height of isosceles triangle $ABD$ and get two equations by the \href{https://artofproblemsolving.com/wiki/index.php?title=Pythagorean\_Theorem}{Pythagorean Theorem}.First, $h^2 + k^2 = 86^2$. Second, $h^2 + (k + m)^2 = 97^2$.Subtracting these two equations, we get $2km + m^2 = 97^2 - 86^2 = (97 - 86)(97 + 86) = 2013$.We then add $k^2$ to both sides to get $k^2 + 2km + m^2 = 2013 + k^2$.We then complete the square to get $(k + m)^2 = 2013 + k^2$. Because $k$ and $m$ are both integers, we get that $2013 + k^2$ is a square number. Simple guess and check reveals that $k = 14$.Because $k$ equals $14$, therefore $m = 33$. We want $\overline{BC} = 2k + m$, so we get that $\overline{BC} = \boxed{(D)~61}$.

 $\phantom{solution and diagram by bobjoe123}$

 



\textbf{Solution 4}

Let $E$ be the foot of the altitude from $A$ to $BX.$ Since $\triangle ABX$ is isosceles $AX=AB=86,EB=EX,$ and the answer is $EC+EB=EC+EX.$ $(EC+EX)(EC-EX)=EC^2-EX^2=(97^2-AE^2)-(86^2-AE^2)=97^2-86^2=2013$ by the Pythagorean Theorem. Only $EC+EX=\boxed{(D)~61}$ is a factor of $2013$ such that $97>EC+EX>EC-EX=\frac{2013}{EC+EX}.$

 



\textbf{Solution 5}

We can apply Stewarts to $\triangle{ABC}$ with cevian $AX.$ If we let $BX = b$ and $CX = c,$ We can write the equation,
 \[bc(b+c) + 86^2(b+c) = 97^2(b) + 86^2(c).\]Now we do a clever set of changes to factor,
 \[bc(b+c) = 97^2(b) + 86^2(c) - 86^2(b+c) = 97^2(b)+86^2(c-(b+c)).\]
 \[bc(b+c)=97^2(b)+86^2(-b) = b(97^2-86^2).\]
 \[c(b+c) = 97^2 - 86^2 = 11\cdot 183 = 3 \cdot 11 \cdot 61.\]Now we can just use the factors of $3 \cdot 11 \cdot 61$ to find what $c$ is and an allowed value of $b+c.$ Using triangle inequality we have 
 \[11 < BC = b+c < 183.\] Therefore applying these restrictions the only value of $b+c$ that works is $\boxed{61}$ with corresponding values of $c=33$ and $b = 28.$

 



\textbf{Solution 6 (Law of Cosines)}

Let $x = BC$, $y = XC$, and $\theta = \angle BCA = \angle XCA$. Using law of cosines on $\triangle ABC$ and $\triangle AXC$ we get
 \[\cos(\theta) = \frac{97^2-86^2 + x^2}{2\cdot97x} = \frac{97^2-86^2 + y^2}{2\cdot97y}\]
 \[(2013 + x^2)y = (2013 + y^2)x\]
 \[2013(x-y) = xy(x-y)\]Because $x-y = \overline{BX}$, we know $x-y$ is nonzero, so
 \[xy = 2013\]Now you can list factors of 2013 and use the triangle innequlity as above solutions to get $x = \boxed{(D)~61}$

 


 


	\item (\textit{AMC 10 10A 2013}) Central High School is competing against Northern High School in a backgammon match. Each school has three players, and the contest rules require that each player play two games against each of the other school's players. The match takes place in six rounds, with three games played simultaneously in each round. In how many different ways can the match be scheduled?


 $\textbf{(A)}\ 540\qquad\textbf{(B)}\ 600\qquad\textbf{(C)}\ 720\qquad\textbf{(D)}\ 810\qquad\textbf{(E)}\ 900$


 
 \textbf{Answer: }E



\textbf{Solution 1}

Let us label the players of the first team $A$, $B$, and $C$, and those of the second team, $X$, $Y$, and $Z$. 

 $\textbf{1}$. One way of scheduling all six distinct rounds could be:

 Round 1: $AX$ $BY$ $CZ$

 Round 2: $AX$ $BZ$ $CY$

 Round 3: $AY$ $BX$ $CZ$

 Round 4: $AY$ $BZ$ $CX$

 Round 5: $AZ$ $BX$ $CY$

 Round 6: $AZ$ $BY$ $CX$

 
The above mentioned schedule ensures that each player of one team plays twice with each player from another team. Now you can generate a completely new schedule by permutating those $6$ rounds and that can be done in $6!=720$ ways.

 $\textbf{2}$. One can also make the schedule in such a way that two rounds are repeated.  

 (a)

 Round 1: $AX$ $BZ$ $CY$

 Round 2: $AX$ $BZ$ $CY$

 Round 3: $AY$ $BX$ $CZ$

 Round 4: $AY$ $BX$ $CZ$

 Round 5: $AZ$ $BY$ $CX$

 Round 6: $AZ$ $BY$ $CX$

 (b)

 Round 1: $AX$ $BY$ $CZ$

 Round 2: $AX$ $BY$ $CZ$

 Round 3: $AY$ $BZ$ $CX$

 Round 4: $AY$ $BZ$ $CX$

 Round 5: $AZ$ $BX$ $CY$

 Round 6: $AZ$ $BX$ $CY$

 As mentioned earlier any permutation of (a) and (b) will also give us a new schedule. For both (a) and (b) the number of permutations are$\frac{6!}{2!2!2!}$ = $90$

 
So the total number of schedules is $720+90+90$ =$\boxed{\textbf{(E)} 900}$.

 



\textbf{Solution 2}

Label the players of the first team $A$, $B$, and $C$, and those of the second team, $X$, $Y$, and $Z$. We can start by assigning an opponent to person $A$ for all $6$ games. Since $A$ has to play each of $X$, $Y$, and $Z$ twice, there are $\frac{6!}{2!2!2!} = 90$ ways to do this. We can assume that the opponents for $A$ in the $6$ rounds are $X$, $X$, $Y$, $Y$, $Z$, $Z$ and multiply by $90$ afterwards.

 
Notice that for every valid assignment of the opponents of $A$ and $B$, there is only $1$ valid assignment of opponents for $C$. More specifically, the opponents for $C$ are the leftover opponents after the opponents for $A$ and $B$ are chosen in each round. Therefore, all we have to do is assign the opponents for $B$. This is the same as finding the number of permutations of $X$, $Y$, and $Z$ that do not have a $X$ in the first two spots, an $Y$ in the next two spots, and a $Z$ in the final two spots.

 


 We can use casework to find this by using the fact that after we put down the $X$'s and $Y$'s first there is $1$ way to put down the $Z$'s (the two remaining spots).

 If $X$'s are put in the middle $2$ spots, then there is $1$ way to assign spots to $Y$, namely the last $2$ spots. (If one of the last two spots are left empty, there will have to be a $Z$ there, which which not valid).

 If $X$'s are put in the last $2$ spots, then there is $1$ way to assign spots to $Y$.

 Finally, if one $X$ was put in on of the middle two spots and one $X$ was put in one of the last two spots, there are $2\cdot2$ ways to assign spots to $X$ and $2\cdot1$ ways to assign spots to $Y$ (one of the first two spots and the remaining spot in the last $2$).

 There are $1+1+2\cdot2\cdot2\cdot1 = 10$ ways to assign opponents to $B$ and $90\cdot10 = 900$ ways to order the games. $\boxed{\textbf{(E)} 900}$

 



\textbf{Solution 3(similar to 2)}

We only need to worry about what the other member of the pair would be for our own competitors. If our competitors are $A,B,C$ and the others are $1,2,3$, then we only need to arrange the numbers $1,1,2,2,3,3$ such that we don't have the same number in the same spots. For our first row of numbers, or our first player, they can be arranged in any of 
 \[\frac{6!}{2!2!2!}\] ways as nothing has happened yet. This contributes $90$ to our final product. Once we determine our second row, or player, the third will immediately follow thus we only need to figure out the second one. WLOG, our first arrangement was $112233$. For our $1$'s, we can put them anywhere that is not in our first $2$ spots. If we put both of them in the same "region", like the $2$'s or the $3$'s, then it will give a different case than if they were in different regions. 

 Case 1 : They are in the same region

 There are $2$ ways that the $1$'s can be in the same region. If, let's say its in the $2$'s region, we try placing the $2's$, we can only place them in the $3$ region in order to fit the $3$'s. Therefore, this entire case produces $2$ arrangements. \\

 Case 2: They are in a different region 

 There are $4$ ways that the $1$'s can be in different regions. For the $2$'s and $3$'s, they will have to go in each other's region with $1$ and the other can be arranged in $2$ ways in the $1$ region giving a total of $8$ for this case.\\

 Finally, $90(2+8) = 900$

 \textbf{Graph Theory Insight} This is a bipartite graph perfect matching problem.  This bipartite graph has $3!=6$ perfect matchings.  Arranging these 6 perfect matchings in 6 rounds gives $6!=720$ ways of scheduling, as Solution 1 states. But Solution 1 did not write about why there are 90 ways to make the schedule that two rounds are repeated, and why 90 is added twice.Let's call round 1 and 2  $P$, round 3 and 4  $Q$, round 5 and 6  $R$.

 \begin{verbatim}
Round 1---->$AX$ $BZ$ $CY$    $P$Round 2---->$AX$ $BZ$ $CY$    $P$Round 3---->$AY$ $BX$ $CZ$    $Q$Round 4---->$AY$ $BX$ $CZ$    $Q$Round 5---->$AZ$ $BY$ $CX$    $R$Round 6---->$AZ$ $BY$ $CX$    $R$
\end{verbatim}

And now the schedule is simply the permutation of $PPQQRR$, which is:$\frac{6!}{2!2!2!}$ = $90$So now we have to figure out how many $PPQQRR$s there are. We first put in the $AX$, $AY$, and $AZ$:$P$---->$AX$

 \begin{verbatim}
$Q$---->$AY$$R$---->$AZ$
\end{verbatim}

Then there are 2 ways to put in $B$'s match in $P$:

 a)

 \begin{verbatim}
$P$---->$AX$ $BY$$Q$---->$AY$$R$---->$AZ$
\end{verbatim}

b)

 \begin{verbatim}
$P$---->$AX$ $BZ$$Q$---->$AY$$R$---->$AZ$
\end{verbatim}

Now we fill in $B$'s match for $Q$ and $R$. Notice how that for each a) and b) there is only one way to fill in $B$'s match for $Q$ and $R$. For a), if we put $BX$ in $Q$, then $BZ$ must be in $R$. But $AZ$ is already in $R$, which makes it impossible. For b), if we put $BY$ in $Q$, $BX$ must be in $R$. Then, $CY$ has to be in both $P$ and $R$ which is impossible. So there is only one way to fill in $B$'s match for $Q$ and $R$.

 a)

 \begin{verbatim}
$P$---->$AX$ $BY$$Q$---->$AY$ $BZ$$R$---->$AZ$ $BX$
\end{verbatim}

b)

 \begin{verbatim}
$P$---->$AX$ $BZ$$Q$---->$AY$ $BX$$R$---->$AZ$ $BY$
\end{verbatim}

Now we fill in $C$'s matchs:

 a)

 \begin{verbatim}
$P$---->$AX$ $BY$ $CZ$$Q$---->$AY$ $BZ$ $CX$$R$---->$AZ$ $BX$ $CY$
\end{verbatim}

b)

 \begin{verbatim}
$P$---->$AX$ $BZ$ $CY$$Q$---->$AY$ $BX$ $CZ$$R$---->$AZ$ $BY$ $CX$
\end{verbatim}

There are 2 $PPQQRR$s, therefore 90 must be added twice.

 


	\item (\textit{AMC 10 10A 2013}) All 20 diagonals are drawn in a regular octagon. At how many distinct points in the interior
of the octagon (not on the boundary) do two or more diagonals intersect?


 $\textbf{(A)}\ 49\qquad\textbf{(B)}\ 65\qquad\textbf{(C)}\ 70\qquad\textbf{(D)}\ 96\qquad\textbf{(E)}\ 128$


 
 \textbf{Answer: }A



\textbf{Solution 1 (Drawing)}

If you draw a clear diagram like the one below, it is easy to see that there are $\boxed{\textbf{(A) }49}$ points.

 
                    \begin{figure}[H]
                        \centering
                        \includegraphics[width=0.5\textwidth]{Solution/AMC_10_10A_2013_25_Solution_Figure_1.png}
                    \end{figure}
                    

 



\textbf{Solution 2 (Working Backwards)}

Let the number of intersections be $x$.  We know that $x\le \dbinom{8}{4} = 70$, as every $4$ vertices on the octagon forms a quadrilateral with intersecting diagonals which is an intersection point.  However, four diagonals intersect in the center, so we need to subtract $\dbinom{4}{2} -1 = 5$ from this count,  $70-5 = 65$.  Note that diagonals like $\overline{AD}$, $\overline{CG}$, and $\overline{BE}$ all intersect at the same point.  There are $8$ of this type with three diagonals intersecting at the same point, so we need to subtract $2$ of the $\dbinom{3}{2}$ (one is kept as the actual intersection). In the end, we obtain $65 - 16 = \boxed{\textbf{(A) }49}$.

 



\textbf{Solution 3 (Answer choices and reasoning)}

We know that the amount of intersection points is at most $\dbinom{8}{4} = 70$, as in solution $2$. There's probably going to be more than $5$ intersections counted multiple times (to get $\textbf{(B) }65$), leading us to the only reasonable answer, $\boxed{\textbf{(A) }49}$.-Lcz

 Note: You can easily prove this by looking at the simple case of the diagonals intersecting in the middle of the octagon.  $4$ major diagonals intersect here and only $1$ intersection point is counted so you can subtract $3$ from $70$.  Then look to the middle area of the octagon.  In this area, if we label the major diagonal as the one where there are $3$ points between the two points forming the diagonal, and the semi-minor diagonal the diagonal where there is one less point between the two diagonal forming points, there are $8$ intersection points of a major diagonal and $2$ semi-minor diagonals.  This means that these eight points would be, not double-counted -which the calculation by Lcz accounts for- but triple-counted.  Thus, taking away one for each of these points is more than enough to see the value of the answer has to be less than or equal to $59$.  Choice A is the only answer that works.

 



\textbf{Solution 4 (Drawing but easier)}

Like solution one, we may draw. Except note that the octagon has eight regions, and each region has an equal number of points, so drawing only one of the eight regions and the intersection points suffices. One of the eight regions contains $8$ points (not including the octagon center). However each adjacent region share one side in common and that side contains $2$ intersection points, so in actuality there are $8 - 2 = 6$ points per region. We multiply this by $8$ to get $6\cdot 8 = 48$ and add the one center point to get $48 + 1 = \boxed{\textbf{(A) }49}$.

 



\textbf{Solution 5 (Case Work with Drawing)}


                    \begin{figure}[H]
                        \centering
                        \includegraphics[width=0.5\textwidth]{Solution/AMC_10_10A_2013_25_Solution_Figure_2.png}
                    \end{figure}
                    

 This problem is a counting problem of combinatoric geometry. There are 2 cases for the above diagram:

 Case 1: Red Dots

 The red dots are the intersection of 3 or more lines. It consists of 8 dots that make up an octagon and 1 dot in the center. Hence, there are $9$ red dots.

 Case 2: Blue Dots

 The blue dots are the intersection of 2 lines. Each vertex of the octagon has 2 purple lines, 2 green lines, and 1 orange line coming out of it. There are 5 dots of intersection on a purple line, 6 dots on a green line, and 5 dots on an orange line. There are $2 \cdot 5+2 \cdot 6+5=27$ dots that come out of 1 vertex, which includes 7 red dots already counted. So there are $27-7 = 20$ blue dots coming out of 1 vertex. There are 8 vertices, but each blue dot is the intersection of 2 lines, corresponding to $2 \cdot 2 = 4$ vertices. So there are $\frac{20 \cdot 8}{4} = 40$ blue dots.

 The number of intersection dots are the sum of the number of red and blue dots. Hence, the answer is $40 + 9 = \boxed{\textbf{(A) }49}$.

 



\textbf{Solution 6 (PIE)}

Notice that there can be a maximum of 4 lines intersecting at any point inside the octagon. Now, let $P_2$ be the number of points where at least $2$ lines intersect, and define $P_3$ and $P_4$ similarly. By PIE, we have that the desired number of points is
 \[P_2 - 2P_3 + 3P_4\]We now separately count the $P_i's$. 

 
Case 1 (P\_2):Notice that if we pick any $4$ points, there is exactly $1$ way $2$ lines can be drawn such that they intersect inside the convex hull of the $4$ points. Thus, $P_2 = {8 \choose 4} = 70$.

 
Case 2 (P\_3):The only ways that $3$ diagonals concur is if they are all the longest (main) diagonals in the octagon, or one of them is a main diagonal and the other $2$ intersect on said diagonal. If they are all the main diagonals, there are $4$ ways. If only one of them is a main diagonal, then there are $4(2) = 8$ ways. In total, there are $8 + 4 = 12$ such points.

 
Case 3 (P\_4): The only way this can happen is if all $4$ diagonals intersect in the center (they are all main diagonals), giving $1$ way.

 
Thus, in total we have
 \[70 - 2(12) + 3 = \boxed{\textbf{(A)} 49}\]such points.

 


\end{enumerate}